\subsection{buffer}\label{buffer}
Das \verb|buffer| Modul stellt einen dynamisch allozierten Datenpuffer
bereit, der einen zusammenhängenden Speicherbereich verwaltet und dessen
Belegung mitführt. Es werden zwei Ausprägungen unterschieden
(\verb|b_type_e|, Listing \ref{buffer_enums}):
\begin{description}
\item[\texttt{LINEAR}] Ein linearer Puffer; Daten liegen ab dem Anfang des
  Speichers, Schreiben beginnt stets bei Index $0$.
\item[\texttt{RING}] Ein Ringpuffer; Schreib\- und Lesepositionen laufen
  über das Speicherende hinaus zyklisch um (\emph{wrap\-around}).
\end{description}
Der Belegungszustand wird über \verb|state_e| ausgedrückt:
\verb|BUFFER_READY| (frei/leer) bzw. \verb|BUFFER_USED| (enthält Daten,
also belegt).

\begin{listing}
\begin{minted}{C}
typedef enum {
    BUFFER_READY, // Buffer frei/leer
    BUFFER_USED,  // enthaelt Daten (belegt)
    BUFFER_CNT
} state_e;

typedef enum {
    LINEAR, // linearer Puffer
    RING,   // Ringpuffer
} b_type_e;
\end{minted}
\caption{Aufzählungen \texttt{state\_e} und \texttt{b\_type\_e}}
\label{buffer_enums}
\end{listing}

\subsubsection*{Zustand}
Ein Puffer wird durch \verb|buffer_t| (Listing \ref{buffert})
beschrieben. \verb|mem| zeigt auf den allozierten Speicher der Größe
\verb|size|; \verb|used| gibt die Anzahl belegter Bytes an und
\verb|first| (für \verb|RING|) den Index des ersten genutzten Bytes. Das
Feld \verb|lbl| (\verb|clabel_u|, siehe Abschnitt \ref{common}) nimmt die
ersten bis zu \verb|CMD_LEN|$-1$ Zeichen des Inhalts als kurzes Label auf,
\verb|id| dient der Identifikation innerhalb eines Pools (Abschnitt
\ref{buffer_pool}).

\begin{listing}
\begin{minted}{C}
typedef struct buffer_s {
    state_e state;
    uint8_t *pl;   // Zeiger auf erstes genutztes Byte
    uint8_t *mem;  // Start des Speichers
    clabel_u lbl;  // erste Zeichen des Inhalts als Label
    int16_t first; // Index des ersten genutzten Bytes
    int16_t used;  // Anzahl belegter Bytes
    int16_t size;  // Groesse des Speichers
    uint8_t id;    // Id des Buffers (im Pool)
    b_type_e type; // LINEAR oder RING
} buffer_t;
\end{minted}
\caption{Definition der \texttt{buffer\_t} Struktur}
\label{buffert}
\end{listing}

\subsubsection*{Lebenszyklus}
{\sloppy\emergencystretch=3em
\begin{description}
\item[\texttt{buffer\_new(size, type)}] Alloziert eine \verb|buffer_t|\-Struktur
  samt Speicher der Größe \verb|size|. Rückgabe: \verb|buffer_t *| (Zeiger auf
  den Puffer, bzw.\ \verb|NULL| bei \verb|size == 0| oder fehlgeschlagener
  Allokation).
\item[\texttt{buffer\_new\_buffer\_t(buffer)}] Re\-initialisiert eine bereits
  bestehende Struktur: Größe und Typ bleiben erhalten, der Speicher wird neu
  angelegt. Rückgabe: \verb|buffer_t *| (\verb|NULL| bei Fehler).
\item[\texttt{buffer\_free(buffer)}] Gibt Speicher und Struktur frei. Rückgabe:
  \verb|em_msg| (\verb|EM_ERR|, wenn der Puffer noch als \verb|BUFFER_USED|
  markiert ist).
\end{description}
\par}

\subsubsection*{Daten schreiben und lesen}
{\sloppy\emergencystretch=3em
\begin{description}
\item[\texttt{buffer\_set(buffer, data, size)}] Schreibt \verb|size| Bytes in
  den Puffer (\verb|LINEAR|: ab Index $0$; \verb|RING|: ab der Schreibposition
  mit Umlauf). Schlägt fehl, wenn \verb|size| den freien Platz übersteigt;
  setzt \verb|state = BUFFER_USED| und aktualisiert das Label. Rückgabe:
  \verb|em_msg| (\verb|EM_OK| bei Erfolg, \verb|EM_ERR| bei ungültigem Puffer
  oder zu wenig Platz).
\item[\texttt{buffer\_get(buffer, data, size)}] Liest bis zu \verb|*size| Bytes
  (begrenzt auf die belegte Menge), entrollt beim \verb|RING| den Umlauf und
  schreibt die tatsächliche Länge nach \verb|*size| zurück. Ist der Puffer
  danach leer, wird \verb|state = BUFFER_READY|. Rückgabe: \verb|em_msg|
  (\verb|EM_OK| bei Erfolg, \verb|EM_ERR| bei ungültigem Puffer oder
  \verb|*size <= 0|).
\item[\texttt{buffer\_transfer(from, to)}] Überträgt den Inhalt von \verb|from|
  nach \verb|to| und unterstützt alle \verb|LINEAR|/\verb|RING|\-Kombinationen.
  Die Kapazität wird \emph{vor} dem Lesen geprüft (transaktional): passt der
  Inhalt nicht, bleibt \verb|from| unangetastet; andernfalls wird \verb|from|
  geleert. Rückgabe: \verb|int16_t| -- die übertragene Länge, bzw.\
  \verb|EM_ERR| bei zu kleinem Ziel.
\end{description}
\par}

\subsubsection*{Verwaltung und Abfrage}
{\sloppy\emergencystretch=3em
\begin{description}
\item[\texttt{buffer\_reset(buffer)}] Setzt Inhalt und Zustand vollständig
  zurück (nullt den Speicher). Rückgabe: \verb|em_msg| (\verb|EM_ERR| bei
  ungültigem Puffer, sonst \verb|EM_OK|).
\item[\texttt{buffer\_clear(buffer)}] Markiert den Inhalt als verbraucht
  (\verb|first += used|, \verb|used = 0|, \verb|state = BUFFER_READY|), ohne
  den Speicher zu nullen. Rückgabe: \verb|em_msg| (\verb|EM_ERR| bei
  ungültigem Puffer, sonst \verb|EM_OK|).
\item[\texttt{buffer\_tolower(buffer)}] Wandelt im belegten Bereich alle
  Großbuchstaben \texttt{A}--\texttt{Z} in Kleinbuchstaben um. Rückgabe:
  \verb|em_msg| (\verb|EM_ERR| bei Nullzeiger, sonst \verb|EM_OK|).
\item[\texttt{buffer\_writeable(buffer)}] Liefert den freien Platz
  (\verb|size - used|). Rückgabe: \verb|int16_t|.
\item[\texttt{buffer\_used(buffer)}] Liefert die Anzahl belegter Bytes.
  Rückgabe: \verb|int16_t|.
\item[\texttt{buffer\_is\_used(buffer)}] Liefert \verb|true|, wenn der Zustand
  \verb|BUFFER_USED| ist. Rückgabe: \verb|bool|.
\item[\texttt{buffer\_get\_clabel(buffer)}] Liefert einen Zeiger auf das Label
  des Puffers. Rückgabe: \verb|clabel_u *|.
\item[\texttt{buffer\_get\_till\_end(buffer)}] Liefert einen temporären
  (statischen) \verb|buffer_t|, der den Bereich ab \verb|first| bis zum
  Speicherende beschreibt -- nützlich für \verb|RING|\-Zugriffe ohne Umlauf.
  Rückgabe: \verb|buffer_t *|.
\item[\texttt{buffer\_check(buffer, reduced)}] Prüft die Struktur auf
  Gültigkeit (Nullzeiger, \verb|size == 0|, und -- außer bei \verb|reduced| --
  \verb|mem == NULL|). Rückgabe: \verb|em_msg| (\verb|EM_OK| bei gültiger
  Struktur, sonst \verb|EM_ERR|).
\item[\texttt{buffer\_print(buffer, title)}] Gibt die Felder und den Inhalt zu
  Debug\-Zwecken aus. Rückgabe: \verb|void|.
\end{description}
\par}
